\begin{pastexam}[label={px:ncku-108-ee-ds-03}]{成大電機所 108 學年度 問答題第 2 題}
（問答）

Given an empty hash table of size 7 and bucket capacity 1 that uses open addressing. The following
sequence of keys is to be inserted: (8, 17, 24, 3, 15, 5). Insert these keys using the following
hash function and approach. Show the resulting hash table step by step. If overflow occurs,
indicate the element that causes the overflow. (6 points)

\qquad h(x) = x \% 7; quadratic probing.
\end{pastexam}

\begin{pastsoln}
二次探測用資結書第~12.3~節的公式：第 $i$ 次試 $(h(x) + i^2) \bmod 7$，$i = 0, 1, 2, \ldots$，也就是從起始
位置跳 $0, 1, 4, 9, \ldots$ 格。 每格（bucket）只能放 $1$ 個鍵。

\textbf{用語}：題目的 bucket／overflow 是 Horowitz 課本的說法。 \textbf{碰撞}（collision）是兩個不同的鍵
算出同一個 bucket；\textbf{溢位}（overflow）是新的鍵算到的 bucket 已經滿了。 bucket 容量是 $1$ 時，起始
bucket 有鍵就滿了，所以碰撞一定同時是溢位，要用探測找別的位置。

\begin{center}
\small
\begin{tabular}{cccl*{7}{c}}
\toprule
鍵 & $h(x)$ & 試過的位置 & 結果 & 0 & 1 & 2 & 3 & 4 & 5 & 6 \\
\midrule
8 & 1 & 1 & 放在 1 & & 8 & & & & & \\
17 & 3 & 3 & 放在 3 & & 8 & & 17 & & & \\
24 & 3 & 3, 4 & \textbf{溢位}，放在 4 & & 8 & & 17 & 24 & & \\
3 & 3 & 3, 4, 0 & \textbf{溢位}，放在 0 & 3 & 8 & & 17 & 24 & & \\
15 & 1 & 1, 2 & \textbf{溢位}，放在 2 & 3 & 8 & 15 & 17 & 24 & & \\
5 & 5 & 5 & 放在 5 & 3 & 8 & 15 & 17 & 24 & 5 & \\
\bottomrule
\end{tabular}
\end{center}
細節：$24 \bmod 7 = 3$，位置 $3$ 被 $17$ 佔了，試 $(3 + 1) \bmod 7 = 4$。 $3 \bmod 7 = 3$，試 $3$、$4$ 都滿，
再試 $(3 + 4) \bmod 7 = 0$。 $15 \bmod 7 = 1$，位置 $1$ 被 $8$ 佔了，試 $2$。 $5 \bmod 7 = 5$，空的。

最後的雜湊表：
\begin{center}
\begin{tabular}{c|ccccccc}
位置 & 0 & 1 & 2 & 3 & 4 & 5 & 6 \\
\hline
鍵 & 3 & 8 & 15 & 17 & 24 & 5 & （空） \\
\end{tabular}
\end{center}

\textbf{溢位的鍵}：$24$、$3$、$15$（起始 bucket 已滿，靠二次探測解決，最後都放進去了）。 沒有任何鍵
「放不進表」：$7$ 是質數，表裡的鍵少於 $(7+1)/2 = 4$ 個之前，二次探測保證找得到空格（資結書
定理~12.3）；$15$ 插入時表裡已有 $4$ 個鍵，但它第二次就找到空格。 如果題目的「overflow」指的是「表裡找不到
位置放」，答案就是沒有發生。
\end{pastsoln}
